\chapter{Умовний оператор if}
\justifying
За замовчуванням хід виконання програми є лінійним: інструкції виконуються послідовно — рядок за рядком, від першого до останнього. Послідовне виконання є базою, але реальні алгоритми вимагають гнучкості та адаптивності.

У більшості задач програма повинна приймати рішення залежно від вхідних даних чи поточного стану обчислень. Для цього використовуються розгалуження (або умовні переходи).

\begin{defn}{Оператор if в Python — це}
	умовний оператор, який дозволяє програмі приймати рішення та виконувати певний блок коду лише тоді, коли виконується задана умова
\end{defn}

В основі прийняття рішень у програмі лежать оператори порівняння (відношення). Вони дозволяють зіставляти значення виразів і повертають результатом булеве значення — \textbf{ True (істина)} або \textbf{False (хибність)}.
\section {Оператори порівнювання}
Python підтримує базові математичні операції відношення:

\begin{pydef}
Дорівнює:                a == b
Не дорівнює:             a != b
Менше ніж:               a < b
Менше або дорівнює:     a <= b
Більше ніж:              a > b
Більше або дорівнює:    a >= b
\end{pydef}

Зверніть увагу:

Оператор перевірки на \textbf{ рівність} позначено двома знаками \textbf{==}, оскільки один знак = в Python зарезервовано для операції \textbf{присвоєння}.

На відміну від багатьох інших мов програмування, Python дозволяє використовувати ланцюгові порівняння (подвійні нерівності), що звично для математиків, наприклад:\textbf{ a < x <= b} замість запису через логічне «і».

Отримані булеві вирази є фундаментом для керування потоком виконання програми: вони використовуються в умовних конструкціях (if, elif, else) та циклах (while).

\section{Синтаксичні форми умовного оператора}

Залежно від кількості напрямків (гілок) виконання програми, в Python застосовують три основні форми розгалуження:

\vspace{1em}

\noindent\textbf{1. Неповна форма (коротке розгалуження)}

Використовується, коли певну дію потрібно виконати \textbf{тільки} за умови, що вираз-умова є істинним (\texttt{True}). Якщо умова хибна (\texttt{False}), блок коду просто пропускається.

\begin{pydef} 
Неповна форма
if вираз_умова1:
    оператор1
    оператор2
\end{pydef}


\vspace{0.5em}

\noindent\textbf{2. Повна форма (двохгілкове розгалуження)}

Застосовується, коли є два альтернативних шляхи: один виконується при істинності умови (\texttt{True}), а інший (блок \texttt{else}) — якщо умова хибна (\texttt{False}).

\begin{pydef}
Повна форма
if вираз_умова1:
    оператор1
    оператор2
else:
    Блок_операторів_3
\end{pydef}

\vspace{0.5em}

\noindent\textbf{3. Багатогілкова форма (багатоальтернативне розгалуження)}

Використовується для перевірки серії послідовних умов за допомогою конструкцій \texttt{elif} (скорочення від \textit{else if}). Умови перевіряються зверху вниз: виконується блок першої умови, яка виявиться істинною, після чого весь оператор \texttt{if} завершує роботу.

\begin{pydef}
Багатогілкова форма
if вираз_умова1:
    Блок_операторів_1
elif вираз_умова2:
    Блок_операторів_2
elif вираз_умова_n:
    Блок_операторів_n
else:
   Блок_операторів_n+1
\end{pydef}

\vspace{0.5em}

\textbf{Важливі синтаксичні особливості Python:}
\begin{itemize}
	\item \textbf{Двокрапка (\texttt{:}):} є обов'язковою в кінці кожного рядка з \texttt{if}, \texttt{elif} та \texttt{else}. Вона позначає початок відповідного блоку коду.
\end{itemize}

\begin{py}{Приклад 1. Користувач вводить число та отримує відповідь яке воно позитивне негативне або нуль.} 
x=int(input("Введіть x "))
if x>0:
    print("x позитивне")
else:
    print("x негативне або нуль")
    input()
\end{py}

\begin{py}{Приклад 2. Користувач вводить два числа та отримує відповідь яке більше або дорівнюють. }
A=input()
B=input()
if A>B: print("А більше") #можна писати у тому ж рядку 1 оп-р
elif A==B: print("А дорівнює В")
else: print("А менше В")
\end{py}

\begin{py}{Приклад 3. Користувач отримує відповідь чи оплачений чек (ціна записана у коді, або можна ввести з клавіатури). (звернути увагу на імена змінних}
print ("Привіт")
price1 = 50 #price1=int(input()) 
price2 = 32 #price2=int(input())
if price1+ price2 > 99 :
    print ("Сума не достатня")
else:
    print ("Чек оплачений")
    print ("До побачення")
\end{py}

\begin{py}{Приклад 4. Вивести на екран місяці відповідної попи року}
s=input("Введіть пору року")
if s=="зима":
    print("грудень")
    print("січень")
    print("лютий")
elif s=="весна":
    print("березень")
    print("квітень")
    print("травень")
elif s=="літо":
    print("червень")
    print("липень")
    print("серпень")
else:
    print("вересень")
    print("жовтень ")
    print("листопад")
print ("До побачення")
\end{py}
\section*{Контрольні питання}

\begin{enumerate}
	\item Що таке умовний оператор \texttt{if} у Python? Для чого він використовується?
	\item Які оператори порівняння підтримує Python? Наведіть приклади.
	\item Чим відрізняється оператор присвоєння \texttt{=} від оператора перевірки на рівність \texttt{==}?
	\item Що таке ланцюгові порівняння у Python? Наведіть приклад.
	\item Які основні синтаксичні форми умовного оператора існують у Python?
	\item У чому полягає відмінність між неповною та повною формою умовного оператора?
	\item Як працює багатогілкова форма умовного оператора з \texttt{elif}?
	\item Яка роль двокрапки \texttt{:} у синтаксисі умовних операторів?
	\item Що відбувається, якщо умова у операторі \texttt{if} є хибною?
	\item Як можна записати умовний оператор у одному рядку?
	
\end{enumerate}
\section*{Практичні завдання}
\begin{enumerate}
	\item Реалізуйте програму, яка за двома введеними числами перевіряє, чи можна знайти частку від ділення першого на друге.
	\item Напишіть програму, яка за введеними координатами точки возначає чверть.
	\item Складіть програму, яка визначає найбільше з трьох введених чисел.
	\item Реалізуйте програму, яка перевіряє, чи введене число є парним або непарним.
	\item Напишіть програму, яка визначає, чи введене число кратне 10.
	\item Складіть програму, яка за введеним віком користувача виводить повідомлення: «дитина», «підліток», «дорослий», «пенсіонер».
	\item Реалізуйте програму, яка перевіряє, чи введений рік є високосним.
	\item Напишіть програму, яка за введеною оцінкою (від 1 до 12) виводить текстове пояснення («незадовільно», «задовільно», «добре», «відмінно»).
\end{enumerate}
